\subsection{Heavy bins and a polynomial wall}\label{w:heavy}

Suppose, towards a contradiction to \Cref{w:weak}, that
\begin{equation}\label{w:eq-counterlevel}
 Y=\{|U|\geq\lambda\},\qquad
 M_Y=\mu(Y)>C_{\mathrm{ind}}mL^{2n}D^{-\sigma}g^{-p}.
\end{equation}
All constants used in the geometric construction below are independent
of $C_{\mathrm{ind}}$.  Translations of the spatial box are immaterial.
For a macro cube with center $z=(x,t)$, call a $1/R$-bin heavy if
\begin{equation}\label{w:eq-heavy}
 \sum_{\alpha\ \mathrm{in\ the\ bin}}|c_\alpha|^2
 \left(1+\frac{|x-y_\alpha-2\omega_\alpha t|}{L}\right)^{-M}
 \geq W,\qquad
 W=X^{-P_0}\frac{mL^{n+1}}{M_Y}.
\end{equation}
The power $P_0$ will be chosen after an inner sparsity exponent $P_1$.

\begin{lemma}[Reduction to broad heavy incidences]\label{w:heavy-reduction}
There is a subset $Y'\subset Y$ with $\mu(Y')\geq cM_Y$ such that,
at each point of $Y'$:
\begin{enumerate}[label=\textup{(\roman*)}]
\item the heavy-bin field has modulus at least $\lambda/2$;
\item the heavy-bin centers are not contained in a strip of any fixed
large constant times $R^{-1}$ about an affine $(n-2)$-plane, nor in a
single quadratic neighborhood of the form in \Cref{w:angular};
\item every heavy bin supplies at least $W/2$ of its weighted
coefficient square mass from trajectories within distance $X^{C_w}L$.
\end{enumerate}
In particular some $n$ heavy bins have direction vectors
$(2\omega,1)$ whose $n$-volume is at least $R^{-C}$ throughout the
bins.
\end{lemma}

\begin{proof}
For each bin restrict $\mu|_Y$ to the macros where that bin is light.
Since $R^4\ll L$, the weight in \eqref{w:eq-heavy} is comparable at
the macro center and at each point of the macro.  Therefore the sum of
its coefficient energies times their packet moments on this restricted
measure is at most $CWM_Y$.

Put $D_1=X^{P_1}$.  Split off the labels whose moment exceeds
$mL^{n+1}/D_1$.  In each bin their total energy is at most
$CD_1X^{-P_0}$; summing over the $O(R^n)$ bins only contributes a fixed
power of $R$.  The other labels satisfy the inner induction hypothesis
at $D_1$, on the restricted measure, with the unchanged original
profiles.  Apply it at threshold $cR^{-n}\lambda$.  Disjoint bin
energies and $p/2>1$ give total exceptional mass at most
\begin{equation}\label{w:eq-light-inner}
 C C_{\mathrm{ind}}mL^{2n}g^{-p}R^{np}D_1^{-\sigma}.
\end{equation}
Choose $P_1$ so that $2D\leq D_1$ and
$un p_{\max}-\sigma P_1+\sigma<0$ with a fixed spare margin.
The later choice of $\kappa$ ensures $D_1\leq L$.  If $D_1$ has
already reached $L^\kappa$, the large-$D$ case is available; otherwise
this is a strict inner induction step.

For each split-off bin, use unrotated cap scaling and the outer
induction at $L/R$.  The ball parameter is $CmR^{n+1}$, the relative
unit density is at most $\min(1,C\chi R)$, and the scaled level
parameter is at least $gR^{-C}$.  Consequently all split-off fields
together contribute at most
\begin{equation}\label{w:eq-light-outer}
 C C_{\mathrm{ind}}mL^{2n}g^{-p}
 R^C(D_1X^{-P_0})^{p/2}.
\end{equation}
The fixed power $R^C$ includes the bin count and threshold splitting;
its exponent can use $p_{\max}$.  Choose $P_0$ sufficiently larger
than $P_1$ so that this expression, and \eqref{w:eq-light-inner}, are
small fractions of the lower bound in \eqref{w:eq-counterlevel} once
$B$ is large.  The thresholds of the $O(R^n)$ light-bin fields sum to
at most $\lambda/2$ when $c$ is small.  This proves \textup{(i)}.

On every adverse macro for \textup{(ii)}, the heavy-bin sum is exactly
a prescribed-bin field of the type in \Cref{w:angular}.  Its mass is
at most the target bound times $CR^{-c_2}D^\sigma$, plus an absolute
negligible error.  Since $\sigma<uc_2$, increasing $B$ makes this a
small fraction.  The construction may choose the adverse hyperplane
and quadratic separately on each macro, as allowed by that lemma.

For \textup{(iii)}, the measure ball bound and the fixed time range
give, by covering tube annuli,
\[
 \int_{|x-y_\alpha-2\omega_\alpha t|>X^{C_w}L}
 \left(1+\frac{|x-y_\alpha-2\omega_\alpha t|}{L}\right)^{-M}
 \,\dd\mu\leq CmL^{n+1}X^{-A_w},
\]
where $A_w$ can exceed $P_0$ by choosing $C_w$ large (the fixed
exponent $M$ was chosen sufficiently large first).  Sum with
$|c_\alpha|^2$ and compare with $WM_Y=X^{-P_0}mL^{n+1}$.
Markov's inequality removes only a small fraction of $Y$ and ensures
the stated contribution in every heavy bin.  Finally, successive
maximal affine heights among the remaining centers give $n$ robustly
independent directions; failure would place all centers in a forbidden
codimension-two strip.  Choosing the strip constant large makes the
volume bound stable throughout the bins.
\end{proof}

Pass to tube coordinates $(x,t)/L$ and replace the measure by
$\nu=\mu/(mL^n)$ in those coordinates.  The relevant domain has diameter
$O(L)$ and
\[
 \nu(B(z,r))\leq Cr^n\quad(r\geq1).
\]
Write $N_Y=\nu(Y)$.  Formula \eqref{w:eq-heavy} becomes
\begin{equation}\label{w:eq-Wnormalized}
 W=X^{-P_0}L/N_Y.
\end{equation}
The incidence tubes have width $X^{C_w}$ in these coordinates.

\begin{lemma}[Incidence mass and a wall]\label{w:wall}
The preceding level set satisfies
\begin{equation}\label{w:eq-incidence-mass}
 N_YW^{n/(n-1)}\leq X^C,
 \qquad N_Y\geq X^{-C}L^n.
\end{equation}
Moreover, a subset of $Y'$ of mass comparable to $N_Y$ lies within
distance $X^C$ of the zero set of a polynomial of degree at most
$X^{C_{\mathrm{part}}}$ in tube coordinates.
\end{lemma}

\begin{proof}
For a robust $n$-tuple, project its directions onto an $n$-dimensional
space where their determinant is at least $R^{-C}$.  Such a space
exists by their $n$-volume lower bound.  A further polynomially fine
angular subdivision makes the projection common to the tuple's
families.  Apply endpoint $n$-linear Kakeya to the projected tubes.

To relate its cube sum to $\nu$, divide into columns of side comparable
to the tube width.  Each tube meets only polynomially many width boxes
over a projected column, uniformly in its label; its projected
direction is bounded away from zero by the determinant condition.
Each box has $\nu$-mass at most $X^C$.  If $a_{i,j}\geq0$ denotes the
incident weight of family $i$ in the $j$th vertical box, H\"older with
$n$ factors followed by the monotonicity of finite sequence norms gives
\[
 \sum_j\prod_{i=1}^n a_{i,j}^{1/(n-1)}
 \leq\prod_{i=1}^n\left(\sum_j a_{i,j}\right)^{1/(n-1)}.
\]
The total coefficient weight in each family is at most one.  The
projected Kakeya bound, its polynomial transversality cost and the
polynomial number of tuples therefore bound the mass seeing weight
$W$ from every family by $X^CW^{-n/(n-1)}$.  This proves the first
inequality in \eqref{w:eq-incidence-mass}.  Substitute
\eqref{w:eq-Wnormalized} to obtain the second.  We also have
$N_Y\lesssim L^n$ from the ball bound.

The separation into cells and a polynomial wall follows the
Fourier-restriction partitioning framework of
Guth \citep[Sections~3.2--3.4]{Guth2016Restriction}.
For completeness, use polynomial partitioning with parameter
$P_2=X^{C'}$ on $\nu|_Y$.  It gives a polynomial of degree $O(P_2)$,
$O(P_2^{n+1})$ sign regions, and mass at most
$CN_Y/P_2^{n+1}$ in each open region.  The finite-set partitioning
theorem follows from iterative polynomial bisection
\citep[Theorem~4.1]{GuthKatz2015}; its integrable-density form is
\citet[Theorem~1.4]{Guth2016Restriction}.  For the present finite Borel
measure, smooth it first and use that bisection argument.  Normalize the
coefficient vector of each of the finitely many factors separately
and take a convergent subsequence.  A compact subset off the limiting
zero set has stable signs; its limiting mass obeys the same bound.
Exhausting an open sign region by such compact subsets proves the
assertion for the original measure.  Mass on the zero set is already
part of the wall.  Connectedness of sign regions is not required.

A central line not contained in the wall meets only $O(P_2)$ open
sign regions.  Let $E_{\mathcal O}$ be the total coefficient energy of
labels whose line visits $\mathcal O$.  Then
$\sum_{\mathcal O}E_{\mathcal O}\lesssim P_2$.  Regions with
$E_{\mathcal O}>C_*P_2^{-n}$ have total mass at most $CN_Y/C_*$ by
the region mass bound.  Choose $C_*$ large.  Away from a width-enlarged
wall, a point and its nearby incident line point are in the same
region.  In every remaining region the preceding projected Kakeya
argument with local energies bounds the off-wall mass of $Y'$ by
\[
 X^CW^{-n/(n-1)}P_2^{-n^2/(n-1)}
 \leq X^{C''}N_YP_2^{-n^2/(n-1)}.
\]
The last inequality uses \eqref{w:eq-Wnormalized} and
$N_Y\lesssim L^n$.  Since
\[
 \frac{n^2}{n-1}-(n+1)=\frac1{n-1}>0,
\]
choosing $C'$ large pays the number of regions.  A fixed positive
fraction of the mass must therefore be near the polynomial wall.
\end{proof}

\subsection{Why the wall is flat}\label{w:curvature}

Testing the second fundamental form on incident line directions is
related to the flat-point argument of
Guth--Katz \citep[Lemma~3.3 and Corollary~3.4]{GuthKatz2010Algebraic}.
Here the incidences are approximate and the ambient dimension is higher;
the quadratic narrow estimate supplies the additional directional
exclusion needed to bound the full Hessian.  We prove the quantitative
reduction to affine plates below.

We now work in macroscopic coordinates, obtained by dividing tube
coordinates by $L$.  The domain is a fixed bounded box, and the wall
and incidence errors are at most
\begin{equation}\label{w:eq-Delta}
 \Delta\leq X^{C_4}/L.
\end{equation}
The exponent $C_4$ is already fixed.  All further enlargements in this
subsection are fixed powers of $X$ and will be small on the
macroscopic scale by the later choice of $\kappa$.

We use the following two elementary consequences of fixed-dimensional
semialgebraic complexity.  Here a description of complexity $X^C$
means that the number and degrees of its defining polynomials are
bounded by such a power; their coefficients are unrestricted.
\begin{enumerate}[label=\textup{(\alph*)}]
\item A bounded set of dimension at most $n-1$ and complexity $X^C$
is covered by $X^{C'}\rho^{-(n-1)}$ boxes of diameter $O(\rho)$.
To see this, use a generic shifted mesh.  A box meeting a component
that does not meet any face contains that component, so the component
bound pays such boxes.  Intersections with grid faces generically lower
dimension, have the same type of complexity bound, and are counted by
induction on dimension, with the zero-dimensional component bound as
base case.
\item A bounded set of complexity $X^C$ has a partition into $X^{C'}$
cells within which any two points can be joined by a path of length
$X^{C'}$.  Use a compatible cylindrical decomposition, including the
enclosing box.  Recursively lift paths in the base to a graph or to
the continuous middle section of a sector, adding vertical segments
inside a sector.  The paths have polynomial description complexity
in fixed dimension.  Each bounded coordinate has polynomially many
monotonicity pieces, so the length is bounded by their total variation.
The paths may be taken continuous and piecewise smooth.
\end{enumerate}
The required cylindrical decomposition, projection and bounded-variable
quantifier-elimination bounds are standard; see \citet{Basu2014}.
These bounds are polynomial in description size when dimension is
fixed and do not depend on coefficient magnitudes.  They control the
number of pieces and path lengths, not graph derivatives.  The latter
are obtained next.

\begin{lemma}[Regular graph pieces and derivative bounds]\label{w:regular-wall}
Let $Z$ be the square-free polynomial wall.  It has a semialgebraic
exceptional subset of dimension at most $n-1$ and complexity $X^C$
such that, outside its $C\rho$-neighborhood, every relevant point of
$Z$ has a local graph coordinate with bounded slope and second and
third graph derivatives bounded by a fixed power of $\rho^{-1}$.
The estimates are independent of the polynomial coefficients.
The mass of the discarded neighborhood in tube coordinates is at most
$X^C\rho L^n$.
\end{lemma}

\begin{proof}
Square-freeness implies that the singular set has dimension at most
$n-1$.  Decompose the bounded regular part into polynomially many
pieces with a specified omitted coordinate $j$, on which the implicit
graph $x_j=\phi(\bar x)$ has slope bounded by a dimensional constant.
Choose a finite set of base directions sufficient to polarize tensors
through order three.  Refine the pieces by the signs, including zero,
of the corresponding fourth directional derivatives of $\phi$.
Implicit differentiation makes these tests semialgebraic with
polynomial complexity.  Put all singular points, lower-dimensional
pieces and frontiers into the exceptional set; use a strictly larger
box so that its artificial boundary is harmless.  The frontier
dimension property gives the stated dimension bound.

At distance $C\rho$ from this set, a straight base segment of length
$c\rho$ lifts in the same graph/sign piece.  Indeed its lift has
bounded speed.  If maximal continuation stopped earlier, its Cauchy
endpoint would be a frontier or a singular point within distance
$\rho$, unless the implicit graph could be continued.  Both alternatives
contradict the stopping assumption.  Global injectivity of the base
projection is unnecessary.

On such a two-sided segment let $g(t)=\phi(u+te)$, so $|g'|\leq C$.
The fixed sign of $g''''$ makes $g'''$ monotone.  If
$|g'''(0)|=A$, on one half-interval of length comparable to $\rho$ its
magnitude is at least $A$, with constant sign.  The midpoint second
difference of $g'$ on that half-interval has magnitude at least
$cA\rho^2$ and at most $4C$.  Thus $A\lesssim\rho^{-2}$.
Apply the same argument on shorter interior intervals, then a
difference quotient for $g'$, to get $|g''|\lesssim\rho^{-1}$.
Polarization bounds the full second and third tensors.

By property \textup{(a)}, the exceptional set is covered by
$X^C\rho^{-(n-1)}$ boxes.  Each becomes a ball of radius $O(\rho L)$
in tube coordinates and has $\nu$-mass at most $C(\rho L)^n$.
Their total mass is $X^C\rho L^n$.  Choose
$\rho=X^{-C_5}$ with $C_5$ large and require $\Delta\ll\rho$.
By \eqref{w:eq-incidence-mass}, this discards an arbitrarily small
fraction of $N_Y$.
\end{proof}

Every remaining point of the level set near the wall can now be
attached to a regular $p\in Z$ at distance $O(\Delta)$, with the
preceding slope and derivative bounds.  No regularity of this
attachment is needed; all bad events below can be described by
existence of an eligible nearby $p$.

\begin{lemma}[Line tests]\label{w:line-tests}
For any prescribed large $A$, choose a fixed $C_6$ sufficiently large.
For every label, outside polynomially many macroscopic time intervals
of total length at most $X^{-A}$, an incidence with any eligible $p$
satisfies
\begin{equation}\label{w:eq-line-tests}
 |v_j-\nabla\phi(p)\cdot\bar v|\leq X^{C_6}/L,
 \qquad |\nabla^2\phi(p)[\bar v,\bar v]|\leq X^{C_6}/L,
 \qquad v=(2\omega,1).
\end{equation}
The exponent bounding the number of intervals is independent of the
sizes of $A,C_6$.
\end{lemma}

\begin{proof}
Lift the projection of the nearby line position onto the local graph.
The lift is $O(\Delta)$ from $p$, and its first and second derivatives
differ from those at $p$ by $O(\Delta X^C)$ by
\Cref{w:regular-wall}.  For the macroscopic line $\ell(s)$ define
\[
 H(s)=\ell_j(s)-\phi(\bar\ell(s)).
\]
At a relevant incidence, $|H|\lesssim\Delta$.
Failure of \eqref{w:eq-line-tests} forces $|H'|$ or $|H''|$ to be at
least $cX^{C_6}/L$, after increasing $C_6$ beyond the fixed derivative
error powers.

Order the regular roots over the projected line parameter by
cylindrical decomposition and refine by these inequalities and signs.
This gives polynomially many smooth branch intervals.  They are graph
sections, not sectors, and their implicit derivatives are the actual
derivatives of $H$.  On an interval where $|H|\lesssim\Delta$ and a
first derivative of fixed sign has magnitude at least $q$, the length
is $O(\Delta/q)$.  For a second derivative of fixed sign and magnitude
at least $q$, a midpoint second-difference estimate bounds the length
by $O((\Delta/q)^{1/2})$.  Set $q=cX^{C_6}/L$ and use
\eqref{w:eq-Delta}.  Increasing $C_6$ makes the sum of lengths at most
$X^{-A}$.  Thresholds change coefficients, not degrees or the number
of quantified variables, so the interval-count exponent does not
increase with $A,C_6$.  Quantifying the existence of any eligible
nearby point preserves this complexity bound and proves the last
assertion as well.
\end{proof}

The bad parameters of one label contribute at most
\begin{equation}\label{w:eq-badline-mass}
 X^C(LX^{-A}+X^C)
\end{equation}
to its enlarged tube's $\nu$-mass: cover the bad intervals in tube
coordinates by width boxes and use the ball bound.  The prefactor and
count exponents are fixed independently of $A,C_6$; the initial
$\Delta$ can be enlarged to include the incidence width.
Choose $A$ large, then $C_6$, and finally let $L$ dominate the fixed
powers.  Sum \eqref{w:eq-badline-mass} with $|c_\alpha|^2$ and compare
with $WN_Y=LX^{-P_0}$.  Markov removes only a small fraction of $Y'$,
and ensures that each heavy bin supplies at least one label passing
both tests.  In fact it supplies positive weighted mass of such
labels.  All assertions are simultaneous at each retained wall point.

\begin{proposition}[Affine plate covering]\label{w:geometry}
A subset $Y_0\subset Y'$ with $\mu(Y_0)\gtrsim M_Y$ is contained
in a union of at most $X^C$ affine hyperplane plates, each of thickness
$X^C$ in tube coordinates.
\end{proposition}

\begin{proof}
At a retained regular point $p$, let $\mathbf n=(\mathbf n_x,
\mathbf n_t)$ be a unit normal.  The first line test gives
\[
 |\mathbf n_x\cdot2\omega+\mathbf n_t|\lesssim X^{C_6}/L
\]
for a supplied bounded frequency, so $|\mathbf n_x|\gtrsim1$.
The exact tangent-frequency hyperplane is consequently well conditioned.

Suppose the graph Hessian has norm greater than $X^{C_7}/L$.
Normalize its quadratic form and restrict it via
$\omega\mapsto\overline{(2\omega,1)}$ to that exact hyperplane.
The resulting polynomial on the hyperplane has coefficient norm
comparable to one.  To verify this, choose coordinates on the tangent
space in which its time-one slice has bounded offset.  Its direction
space and one bounded point span the tangent space with bounded
condition number because $|\mathbf n_x|\gtrsim1$.  A homogeneous
quadratic has the form
\[
 Q(y,t)=y^TCy+2t b\cdot y+at^2;
\]
its restriction to $t=1$ records all of $C,b,a$.  The bounded graph
slope makes base projection another bounded isomorphism.  This proves
the coefficient comparison.  Divide the restricted polynomial by its
coefficient norm, so that it has norm exactly one as required in
\Cref{w:angular}; this changes the following bounds only by fixed factors.

Move each supplied frequency to the exact tangent hyperplane.  Its
distance is $O(X^{C_6}/L)$, and by the second line test its normalized
quadratic value there is at most
\[
 C X^{C_6-C_7}+C X^{C_6}/L.
\]
Choose $C_7$ large relative to $C_6,u$ and then make $L$ dominant.
This quantity is less than $R^{-4}$.  Every heavy-bin center is within
$O(R^{-1})$ of these frequencies, so all heavy bins are trapped in a
forbidden quadratic neighborhood from \Cref{w:heavy-reduction}.
The contradiction proves
\begin{equation}\label{w:eq-smallhess}
 \norm{\nabla^2\phi(p)}\leq X^{C_7}/L.
\end{equation}
This is where the quadratic narrow estimate is essential: tangent
directions on a nonflat ruled quadric may span an affine hyperplane,
but their quadratic values still detect its curvature.

For each fixed graph coordinate $j$, partition the exact semialgebraic
subset defined by regularity, the bounded slope condition and
\eqref{w:eq-smallhess} into cells with property \textup{(b)} above.
Use paths inside these cells, not paths inside the larger regular
wall.  Even if the base projection has several sheets, the implicit
gradient
\[
 g_j(p)=-\partial_{\mathrm{base}}P(p)/\partial_jP(p)
\]
is a single-valued function of the wall point.  Along a piecewise
smooth path $p$ in this subset, with base coordinate $u$,
\[
 \dd g_j=\nabla^2\phi_j\,\dd u,\qquad
 \dd p_j=g_j\cdot\dd u.
\]
For path length $D_0\leq X^C$, these identities bound the gradient
change by $X^{C_7}D_0/L$ and the height error from the starting tangent
plane by $X^{C_7}D_0^2/L$.  A projected loop reaching a different
sheet obeys the same identities.  Thus every cell lies within
$X^C/L$ of one affine plane in macroscopic coordinates.  Add the
$O(\Delta)$ attachment error and return to tube coordinates.  The
number of cells and the plate widths are both fixed powers of $X$,
as claimed.
\end{proof}
